\chapter{Thunderstorm asthma}
\index{Thunderstorm asthma}

\medskip
\noindent\textit{Educational reference; always seek professional advice for individual care.}

\section*{Definition and Overview}
Thunderstorm asthma is a sudden, sometimes epidemic surge of acute asthma symptoms triggered when a thunderstorm interacts with high airborne allergen levels---classically grass pollen. Storm outflows can concentrate pollen grains, rupture them into tiny respirable starch granules, and sweep them down to ground level where people inhale a massive allergenic load deep into the airways. Within minutes to hours, susceptible individuals develop wheeze, chest tightness, cough, and severe breathlessness. The phenomenon gained national attention in many countries after the November 2016 Melbourne event, which caused thousands of emergency presentations and multiple deaths. Thunderstorm asthma can affect people with known asthma and also people with seasonal allergic rhinitis who have never been diagnosed with asthma. This chapter explains the mechanism, who is at risk, how warnings work, and what to do before and during a high-risk storm season.

\section*{Epidemiology}
Epidemic thunderstorm asthma is uncommon but high-impact. Documented events have occurred in many countries (especially Victoria during late spring ryegrass pollen peaks), as well as in the United Kingdom and elsewhere. The 2016 Melbourne epidemic remains one of the largest recorded: health services were overwhelmed in hours. Risk concentrates in people with springtime hay fever, outdoor exposure during the storm front, and those with poorly controlled or undiagnosed asthma. Non-English-speaking communities had disproportionate harm in 2016, highlighting warning and education gaps. Seasonal public health alerts now operate in Victoria and inform practice elsewhere when pollen and storm forecasts align. Climate, land use, and pollen seasons influence year-to-year risk.

\section*{Aetiology and Pathophysiology}
\begin{itemize}
    \item \textbf{High grass pollen load:} ryegrass and related pollens dominate many southern local spring seasons
    \item \textbf{Thunderstorm outflow winds:} downdrafts carry pollen from altitude to breathing zones ahead of rain
    \item \textbf{Pollen rupture:} moisture and electrical conditions fragment grains into particles small enough to reach lower airways
    \item \textbf{IgE-mediated bronchospasm:} sensitised mast cells release mediators causing smooth muscle constriction, oedema, and mucus
    \item \textbf{Rapid onset widespread exposure:} many people are hit simultaneously, creating epidemic demand
    \item \textbf{Not ordinary ``getting wet in rain'':} the allergenic aerosol, not cold rain alone, drives the asthma wave
    \item \textbf{Co-factors:} smoke, pollution, and respiratory viruses can worsen vulnerability in a given season
\end{itemize}

\section*{Clinical Features}
\begin{itemize}
    \item \textbf{Sudden wheeze and chest tightness:} during or shortly after a spring thunderstorm
    \item \textbf{Severe breathlessness and cough:} may escalate over minutes
    \item \textbf{Hay fever flare:} sneezing, itchy eyes, and rhinorrhoea often coexist
    \item \textbf{Night-time clustering:} many events strike late afternoon or evening with storm lines
    \item \textbf{People without prior asthma diagnosis:} first-ever attack is well described
    \item \textbf{Hypoxia and silent chest:} life-threatening severity in a subset
    \item \textbf{Anxiety and panic:} compound the sensation of air hunger
\end{itemize}

\section*{Differential Diagnosis}
\begin{itemize}
    \item \textbf{Usual acute asthma exacerbation:} similar treatment but different population trigger pattern
    \item \textbf{Anaphylaxis:} look for urticaria, throat swelling, hypotension---adrenaline pathway if present
    \item \textbf{Panic attack:} must not exclude asthma; objective wheeze and peak flow help
    \item \textbf{Cardiac asthma / heart failure:} older patients, oedema, ischaemic history
    \item \textbf{Smoke or chemical inhalation:} bushfire or industrial context
    \item \textbf{Foreign body or vocal cord dysfunction:} selected atypical cases
\end{itemize}

\section*{When to Seek Medical Care}
If you have asthma or spring hay fever, see a GP before each pollen season to review inhalers and write an action plan. During a thunderstorm asthma warning, stay indoors with windows closed and follow your plan at the first symptom.

\textbf{Contact your local emergency services for:}
\begin{itemize}
    \item Severe difficulty breathing, inability to speak full sentences, blue lips, or collapse
    \item No improvement after reliever inhaler doses as per action plan
    \item Exhaustion, confusion, or drowsiness with asthma symptoms
\end{itemize}

\section*{Diagnosis and Investigations}
\begin{itemize}
    \item \textbf{Clinical diagnosis in the acute event:} history of storm exposure plus bronchospasm
    \item \textbf{Oxygen saturation and peak expiratory flow:} severity assessment
    \item \textbf{Response to bronchodilators:} supports asthma physiology
    \item \textbf{Allergy assessment later:} skin prick or specific IgE to ryegrass and other pollens
    \item \textbf{Spirometry when stable:} confirms variable airflow obstruction and guides long-term therapy
    \item \textbf{Public health surveillance:} ED presentation spikes define epidemics operationally
\end{itemize}

\section*{Management}
\begin{itemize}
    \item \textbf{Acute attack:} sit upright, take salbutamol via spacer as per plan, repeat as directed, add oxygen and nebulised therapy in emergency care, systemic corticosteroids early in significant attacks
    \item \textbf{Emergency escalation:} ambulance, ED protocols for life-threatening asthma including magnesium and ventilatory support if needed
    \item \textbf{Seasonal preventive therapy:} inhaled corticosteroids for asthma; optimise adherence months before spring if indicated
    \item \textbf{Allergic rhinitis control:} intranasal corticosteroids and antihistamines reduce overall airway irritability
    \item \textbf{Allergen immunotherapy:} specialist option for severe grass pollen allergy in selected patients, planned well outside acute storms
    \item \textbf{Device technique:} spacer use multiplies drug delivery compared with puffer alone for many people
    \item \textbf{Action plans and spare relievers:} critical before high-risk days
    \item \textbf{Hospital surge plans:} community messaging and ambulance pre-alert during forecast events
\end{itemize}

\section*{Complications}
\begin{itemize}
    \item \textbf{Respiratory arrest and death:} documented in epidemic settings
    \item \textbf{ICU admission and intubation:} severe cases
    \item \textbf{Pneumothorax and barotrauma:} rare ventilation complications
    \item \textbf{Ongoing poorly controlled asthma:} after a severe first presentation
    \item \textbf{Health system overload:} delayed care for concurrent emergencies during epidemics
\end{itemize}

\section*{Prognosis and Outcomes}
With rapid bronchodilator and steroid treatment, most people recover. Outcomes are worst when there is delay, no reliever available, undiagnosed asthma, or extreme allergen load. Since 2016, improved warnings, clinician awareness, and public education aim to reduce mortality. Long-term prognosis after a thunderstorm asthma attack is good if maintenance therapy and allergen strategies are put in place.

\section*{Prevention and Self-Care}
\begin{itemize}
    \item Know whether you have spring hay fever or asthma; get a written asthma action plan
    \item Check pollen counts and thunderstorm asthma warnings in season (especially Victoria's alert system)
    \item On high-risk days, stay indoors before the storm front, close windows, and avoid outdoor exercise
    \item Carry a reliever inhaler and spacer; ensure the inhaler is in date and not empty
    \item Take preventer medicines every day as prescribed---not only when wheezy
    \item Treat hay fever properly through spring rather than tolerating constant symptoms
    \item After any severe attack, book GP review within days to adjust treatment
    \item Help family members who do not speak English well to understand warnings and 000 use
\end{itemize}

\section*{Sources and Further Reading}
\begin{itemize}
    \item NHS. \textit{Asthma.} https://www.nhs.uk/
    \item Mayo Clinic. \textit{Asthma attack.} https://www.mayoclinic.org/
    \item MSD Manual. \textit{Asthma.} https://www.msdmanuals.com/
    \item MedlinePlus. \textit{Asthma.} https://medlineplus.gov/
    \item Better Health Channel. \textit{Thunderstorm asthma.} https://www.betterhealth.vic.gov.au/
    \item Healthdirect. \textit{Thunderstorm asthma.} https://www.healthdirect.gov.au/
\end{itemize}
